\section{Forward sampling of a prepared cell}
\label{ap:scalar}

The prepared certificate must first produce terminal inputs to which
the absolute increment applies.  If $Y$ denotes their mean, with
stopped paths counted as zero, our immediate goal is
$\E Y\ge(1-\eta)a$.  This estimate uses one-site sampling and
proportional control of the points removed at each layer.  It does not
use the scalar comparison for scores.

The points removed at active layers are specified by a finite-ring
construction.  We give its full definition and the estimates needed
for descent here.  Its proof appears in Section~\ref{ap:ascent},
where the same construction is used to return the increment: its
affine planes solve the congruence constraints without subdividing
at the full modulus.

We keep the one-layer notation of Section~\ref{ap:paths}. In particular,
the layer is $j$, the parent box is $U$, the chart width is $v$, and
\[
 C(\psi(t))-c=y(t)+\beta(t),\qquad
 m_j=\frac{v^{\dim W}\prod_i k_i}{\covol(W\cap\Z^{d_j})}.
\]
Write $I$ for the inactive sheet axes, including the omitted axes;
the latter satisfy $z_i(t)=0$ identically. The parameters belonging to these
axes form $t_I$; all other parameters form $t_\parallel$. Throughout
this section a true site means a point $u\in U$ in the buffered
uniqueness chart, with its unique integer lift $\beta(u)$, and
$X=(u,\beta(u))$.

The congruence outputs will be denoted by
\begin{equation}\label{ap:scalar:outputs}
 b(t)=(\psi(t),w(t),z_{\rm act}(t)),\qquad
 b(X)=(u,w(X),z_{\rm act}(X)).
\end{equation}
Here $w$ is the deck coordinate, not a width. Their total dimension is
$m_b$, and their respective degree bounds, or tags, are $1,j,j$.
Equal tags are combined when using modular estimates. The congruence
output $b$ in Equation~\eqref{ap:scalar:outputs} is distinct from
the determining block $b_j$.

\subsection{Individual moduli and selected affine planes}

The eventual bookkeeping modulus is
\begin{equation}\label{ap:scalar:large-modulus}
 \mathfrak M=\operatorname{lcm}(1,\ldots,Z)
       =\prod_{\varpi\le Z}\varpi^{a_\varpi},\qquad
 a_\varpi=\left\lfloor\frac{\log Z}{\log\varpi}\right\rfloor.
\end{equation}
The logarithm of $Z$ has a pre bound: warm pre at a passive layer and
cold pre at an active layer. The much larger number $\log\mathfrak M$
is not charged to that bound. No equidistribution assertion below uses
the whole modulus $\mathfrak M$, except when it happens to have an
allowed size. We first forecast the detection, modeling, and symbol
bounds, and only then choose $Z$ to contain their required prime powers.

We shall repeatedly use the following consequence of
Lemmas~\ref{cube:modular-exceptions} and~\ref{cube:modular-decay}.
Increase the structural number of pure-free columns before choosing
any numerical accuracies. Apart from an arbitrarily small prescribed
warm probability of paths, the leading homogeneous parts of
Equation~\eqref{ap:scalar:outputs}, in $t_\parallel$, satisfy the modular
rank hypotheses through $Z$, with exceptional-power product
\begin{equation}\label{ap:scalar:modular-goodness}
              \prod_{\varpi\le Z}\varpi^{b_\varpi}\le D'_0,
              \qquad \log D'_0\text{ warm bounded}.
\end{equation}
Indeed the pure-free coefficients contain the independent disjoint
monomials required by Lemma~\ref{cube:modular-exceptions}.
Lemma~\ref{ap:path:ensemble} transfers their product residue laws at
each individual witness modulus. The exceptional-product proof uses
only such witnesses, so its bound is independent of the number of
tested primes. We call paths satisfying
Equation~\eqref{ap:scalar:modular-goodness} modular-good.

At an active layer we mark exceptional parameter values using the
number of suitable affine sections of the congruence map. The
sections satisfy an exact identity: their $s$ coordinates prescribe
the change in the output, while their $h$ coordinates change the
inactive linear variables. These same sections will be used in the
return argument.

\begin{lemma}[Selected planes and flags]\label{ap:scalar:flags}
Let $x$ be a vector of $v_f$ pure-free columns and let $r_0$ be the
vector of additional coefficient-one linear inactive columns. Fix all
other parameter entries. Put $\ell=m_b+\dim r_0$, and let the $i$th
component of $b$ have degree bound $d_i\in\{1,j\}$. Set
\[
 J=\sum_{i=1}^{m_b}\left\{\binom{\ell+d_i}{d_i}-1\right\},
 \qquad
 n_{\rm der}=\sum_{i=1}^{m_b}\binom{\ell+d_i-1}{d_i-1}.
\]
For every sufficiently small fixed $c_\star>0$, a polynomial structural
choice of $v_f$ has the following properties. At a prime $\varpi$, let
$\epsilon_\varpi=c_\star\varpi^{-2}$. Count full-column-rank matrices
$V\in\operatorname{Mat}_{v_f\times\ell}(\mathbb F_\varpi)$ satisfying
the formal identity
\begin{equation}\label{ap:scalar:plane}
 b\bigl(x+V(s,h),r_0+h\bigr)=b(x,r_0)+s,
       \qquad (s,h)\in\mathbb F_\varpi^{m_b}\times
                           \mathbb F_\varpi^{\dim r_0}.
\end{equation}
There are $J$ nonconstant scalar coefficient equations in this identity.
Select a solution $V$ only if both of the following conditions hold
at every basepoint on its plane: the full number of solutions is at
most $(1+\epsilon_\varpi)\varpi^{v_f\ell-J}$, and the
coefficient-equation Jacobian in $V$ has rank $J$ at the transported
solution.

At depth $a\ge1$, let $w_{\varpi,a}(t)$ be the number of solutions of
Equation~\eqref{ap:scalar:plane} modulo $\varpi^a$ with selected
reduction, divided by $\varpi^{a(v_f\ell-J)}$. Then:
\begin{enumerate}
\item Selection is invariant under changing basepoint along the plane,
and $w_{\varpi,a}(t)=w_{\varpi,1}(t\bmod\varpi)$ exactly.
\item Always $0\le w_{\varpi,a}\le1+\epsilon_\varpi$. Under the uniform
nonconstant pure-free coefficient law, at every fixed basepoint,
\[
 \Pr\bigl(\,|w_{\varpi,1}-1|>C\epsilon_\varpi\,\bigr)
                  \le c_\star\varpi^{-10}.
\]
Here $C$ is fixed, and the exponent $10$ can be increased by enlarging
the structural column count.
\item Define a bad flag at $\varpi$ by the failure of the lower bound
$w_{\varpi,1}\ge1-C\epsilon_\varpi$, and put
\begin{equation}\label{ap:scalar:selected-weight}
 W_*(t)=\prod_{\varpi\le Z}w_{\varpi,a_\varpi}(t).
\end{equation}
Then $W_*\le1+O(c_\star)$ everywhere and $W_*\ge1-O(c_\star)$ when
there is no bad flag. For a squarefree product $d$ of primes, the
uniform-model probability of joint bad flags is at most $d^{-10}$;
the sum of the single-prime probabilities is $O(c_\star)$.
\item In the actual ensemble these estimates transfer, with negligible
errors, at the individual moduli used to test them. After averaging
the center over unrestricted Haar measure, the total one-site loss
from flags, with any nonnegative bounded site insertion, is at most
$O(c_\star)$ times its parent mean, plus an arbitrarily small
cold-pre error.
\end{enumerate}
All flags depend only on nonconstant coefficient data and on
$t\bmod\varpi$, never on the constant coefficient of $b$.
\end{lemma}

The proof of Lemma~\ref{ap:scalar:flags} is given in
Section~\ref{ap:ascent}, after the comparison measure is defined.  For the descent we use its one-site loss bound;
the full count and translation invariance enter the return.

\subsection{Descent of the density certificate}

We pass the density certificate down the stopped tree of
Section~\ref{ap:path:stopped-tree}, using the feasible-center set
\eqref{ap:scalar:feasible-center} and its Haar measure $e_j$.
Once the absolute increment supplies a quantitative score, the
return will use Proposition~\ref{ap:scalar:comparison} in the reverse
order of layers.

The density certificate supplies a strict root tuple. Its layer-$j$
torus point lies in the feasible set with room at least
$\gamma w_j/2$ in each coordinate.
Lemma~\ref{ap:path:chart} applied to a smaller ball about that point
therefore gives a cold-pre bound on $\log(1/e_j)$ at layer $j$.
With $v$ sufficiently small inside the gap, every such path satisfies the enlarged
weight-$j$ constraint. Conversely, any evaluation satisfying its
strict constraint must have a center in
Equation~\eqref{ap:scalar:feasible-center}, with the raw determining
slot specified in the stopped tree.

\begin{proposition}[Forward mass]\label{ap:scalar:forward}
For any sufficiently small fixed $\eta>0$, a prepared certificate at
threshold $a$ gives a surviving set of terminal injective paths whose
expected terminal input mean, with stopped paths counted as zero, is
at least $(1-\eta)a$. The input is pulled back successively and masked
by no bad flag at every active layer. Its support is progression-free.
At least a fraction $\eta a$ of terminal paths have mean at least
$(1-2\eta)a$, after harmless adjustment of the loss constants.

The rank-stop probability at layer $h$ need only have inverse
logarithmic cost polynomial in the data at layers $\ge h$. Discards
immediately after that layer can be requested warm-small after
charging the feasible-center proportions of layers $>h$. Thus this
proposition does not require a lower-layer width in a higher-layer
rank budget.
\end{proposition}

\begin{proof}
Stop a prefix when it fails its current rank test. Also discard
noninjective paths and failures of
Equation~\eqref{ap:scalar:modular-goodness}. Preparation tests the
actual conditioned procedure, including earlier stops. Its requested
rank-stop probabilities are arbitrarily small at cold cost in the
current and higher layers. Modular-goodness discards can be made
arbitrarily warm-small uniformly over feasible centers, with $D'_0$
still warm. Injectivity errors can be requested at the same accuracy
by the later side choice. Choose these errors small enough to pay for
the indicated products of later $e_i$.

Fix a nonstopped, rank-sufficient parent at layer $j$. Average first
over unrestricted Haar centers. The one-site calculation in
Lemma~\ref{ap:path:ensemble} says that, with any bounded nonnegative
site insertion, this average is at least $(1-o(1))$ times the parent
mean off a warm-thickness spatial boundary strip, up to an arbitrarily
small cold-pre error. It also has the corresponding uniform upper
density bound. Apply this first with the input times all remaining
strict constraints. Lemma~\ref{ap:scalar:flags} bounds the bad-flag
loss by a small proportional part of the parent mean and a negligible
additive error. In estimating the boundary loss it suffices to keep
the current strict constraint, since all other factors are at most
one. Choose a smooth torus majorant $G_j$, equal to one on that
strict constraint and supported inside the feasible set
\eqref{ap:scalar:feasible-center}. Then $0\le G_j\le1$ and
$\int G_j\le e_j$. If a spatial boundary cutoff has thickness
$\theta$ and mean $O(n\theta)$ on this parent, the marginal rank
estimate proved with Lemma~\ref{ap:path:grid} gives
\[
 \E_U 1_{\rm boundary}\,1_{\rm current\ strict}
       \le O(n\theta)e_j+\delta,
\]
with an arbitrarily small cold-pre error $\delta$. Smooth spatial
majorants may be used for the boundary indicator. Choose
$\theta$ warm-small relative to $a\prod_{i>j}e_i$, and then choose
$\delta$ cold-small relative to $a\prod_{i\ge j}e_i$.
This makes the boundary contribution sufficiently small without
charging $e_j^{-1}$ to the warm thickness; only later feasible-center
proportions enter after cancelling $e_j$.

Every strict contribution uses a feasible center. Conditioning on
those centers, and then dropping the now processed strict constraint,
shows that the expected new strict-and-input mass is at least
\begin{equation}\label{ap:scalar:strict-step}
       e_j^{-1}(1-\epsilon)\,S
                          -O\left(\epsilon a\prod_{i>j}e_i\right),
\end{equation}
where $S$ is its parent mass and $\epsilon>0$ is an arbitrarily small
fixed step loss. For enlarged constraints the comparison has the
opposite direction. On every feasible-center path the processed
enlarged constraint holds. The unrestricted mixed one-site upper
bound therefore gives parent enlarged mass at least
\begin{equation}\label{ap:scalar:enlarged-step}
          (1-\epsilon)e_j\,T_{\rm child}
                            -O\left(\epsilon\prod_{i\ge j}e_i\right).
\end{equation}
This upper comparison of nonnegative child masses remains valid
after discards. All absolute errors in these statements are requested
before the late rank and side choices, much smaller than the displayed
scales.

Here is the tree bookkeeping explicitly. Let $S_j$ and $T_j$ be the
strict-and-masked-input mass and the enlarged-only mass, respectively,
averaged over prefixes arriving at the rank test at layer $j$, with
the remaining constraints evaluated on their parent. The probability
space retains stopped prefixes as zero mass. Write
\[
      s_j=\frac{S_j}{\prod_{i\ge j}e_i},\qquad
      t_j^*=\frac{T_j}{\prod_{i\ge j}e_i},
\]
with empty products at $D+1$. A rank stop at $h$ loses at most its
probability divided by $\prod_{i\ge h}e_i$ from normalized strict
mass; discarding the output of layer $h$ loses at most its probability
divided by $\prod_{i>h}e_i$. Request the former errors much smaller
than $\epsilon a\prod_{i\ge h}e_i$ and the latter much smaller than
$\epsilon a\prod_{i>h}e_i$. These are exactly the cold and warm
dependencies in the statement. Integration of
Equations~\eqref{ap:scalar:strict-step} and
\eqref{ap:scalar:enlarged-step} now gives
\begin{equation}\label{ap:scalar:mass-recurrences}
 s_{j+1}\ge(1-O(\epsilon))s_j-O(\epsilon a),\qquad
 t_j^*\ge(1-O(\epsilon))t_{j+1}^*-O(\epsilon).
\end{equation}
In the second inequality restrict first to good parents; any later
stopping can only decrease its right-hand nonnegative mass.

At the terminal level $t_{D+1}^*$ is the full survival probability,
which is $1-O(D\epsilon)$ by the requested errors. Thus
Equation~\eqref{ap:scalar:mass-recurrences} implies
\[
 \E B^+=T_1\ge(1-O(D\epsilon))\prod_{i=1}^D e_i.
\]
The certificate $S_1>aT_1$ gives $s_1>a t_1^*$, and the other
recurrence yields
\[
 s_{D+1}=\E[\text{terminal input mean};\text{survival}]
                            \ge(1-O(D\epsilon))a.
\]
Choose $\epsilon$ and $c_\star$ so that the total is at most $\eta$.
If $Y\in[0,1]$ denotes this terminal mean, with $Y=0$ on stopped
paths, then
\[
 \Pr\bigl(Y\ge(1-2\eta)a\bigr)
       \ge \E Y-(1-2\eta)a\ge\eta a.
\]
Finally, an injective integer-affine pullback of a function with
progression-free support again has progression-free support; each
flag mask only removes points. The terminal sides were chosen to
satisfy the dimension-dependent absolute-rule cutoff. That rule
therefore applies on the productive paths at threshold
$(1-2\eta)a$, yielding its uniform multiplicative gain and quantitative
surplus, or its impossibility alternative. All losses here may be
fixed much smaller than that gain.
\end{proof}
